\begin{ouverture}[college3]{15}{Mesurer la tension électrique}{Partie 3 · L’énergie, ses transferts et ses conversions}{Peut-il y avoir une tension sans courant ?}
\objectif{Brancher un voltmètre aux bornes d’un dipôle.}
\objectif{Appliquer additivité en série et unicité en dérivation.}
\objectif{Distinguer tension, intensité et conditions de fonctionnement.}
\prerequis{Nœuds, branches et unités d’intensité.}
\end{ouverture}
\pcPage{Le cours}{Mesurer la tension électrique}
\phantomsection\label{ess-15}
\essSection{1}{La tension entre deux points}
Une tension électrique se définit entre deux points ou aux bornes d’un dipôle. Elle se note $U$ et s’exprime en volts (V). Une source peut présenter une tension entre ses bornes même quand le circuit est ouvert et que le courant extérieur est nul. Tension et intensité sont donc deux grandeurs différentes ; leurs appareils ne se branchent pas de la même manière.\par
\pcFigure{figures/c15-voltmetre.pdf}{Le voltmètre est branché en dérivation aux bornes de la résistance ; l’ampèremètre, s’il existe, reste en série.}
\essSection{2}{Utiliser le voltmètre}
Choisir la fonction tension continue, les bornes V et COM, et un calibre adapté. Relier les pointes aux deux points étudiés sans ouvrir la branche. Pour une lecture positive aux bornes d’un générateur continu, V est reliée à la borne positive et COM à la borne négative. Inverser les pointes inverse le signe de la tension mesurée. On n’utilise ici que des sources scolaires autorisées.\par
\clearpage
\essSection{3}{En série : additionner}
Pour deux récepteurs en série alimentés par un générateur, avec les polarités choisies de façon cohérente, $U_G=U_1+U_2$. Il s’agit de l’additivité des tensions dans la boucle simple étudiée. Les tensions des récepteurs ne sont pas nécessairement égales : elles dépendent de leurs caractéristiques. Diviser la tension du générateur par deux exige des hypothèses supplémentaires.\par
\pcFigure{figures/c15-lois.pdf}{En série, les tensions des récepteurs s’additionnent. En dérivation sur les mêmes nœuds, elles sont égales avec la même polarité.}
\essSection{4}{En dérivation : comparer les mêmes points}
Deux dipôles reliés aux mêmes deux nœuds ont la même tension avec la même orientation : $U_1=U_2=U_G$ dans le montage directement alimenté. Un fil idéal présente une tension nulle à ses bornes. Un interrupteur fermé est modélisé comme un fil ; ouvert dans une boucle simple, il peut supporter une tension non nulle. Les généralisations hors de cette configuration demandent d’examiner le circuit.\par
\begin{pcEncadre}{method}{Un exemple expliqué}Un générateur fournit $6{,}0\unite V$ à deux lampes en série. Si $U_1=2{,}4\unite V$, alors $U_2=6{,}0-2{,}4=3{,}6\unite V$. Si les deux lampes étaient reliées directement en dérivation au générateur idéal de $6{,}0\unite V$, chaque tension serait $6{,}0\unite V$ : c’est un autre montage.\end{pcEncadre}
\pcPage{Expérimenter et retenir}{Mesurer la tension électrique}
\phantomsection\label{activite-15}
\begin{pcEncadre}{method}{Investigation en classe}\textbf{Question.} Comment la tension se répartit-elle ? Avec une source scolaire adaptée et deux dipôles autorisés, réaliser un montage série puis un montage en dérivation. Dessiner avant de mesurer. Relever $U_G$, $U_1$ et $U_2$ avec les mêmes conventions de polarité. Comparer la somme en série et les valeurs en dérivation. \textbf{À rendre :} points de connexion V-COM, unités, résultats et discussion de la précision. Les tensions domestiques sont étudiées uniquement sur document.\end{pcEncadre}
\pcSousTitre{Ce que je dois savoir expliquer}\begin{questions}
\item Une tension se mesure entre deux points, en volts.
\item Le voltmètre se branche en dérivation.
\item Une tension non nulle peut exister dans un circuit ouvert.
\item En série, les tensions des récepteurs s’additionnent dans la boucle simple.
\item En dérivation, mêmes deux nœuds signifie même tension orientée.
\end{questions}
\begin{pcEncadre}{attention}{Une erreur à éviter}Ne pas affirmer « tension identique partout » : c’est l’intensité qui est identique dans une branche série. La tension concerne toujours deux points.\end{pcEncadre}
\pcRepere{cor-activite-15}{Pistes pour l’investigation}\pcGrille
\clearpage\markright{Je démarre}
\rubrique{Je démarre}
\par\noindent\begin{minipage}{\linewidth}
\begin{exercice}[1]{La bonne grandeur}
\begin{questions}
\item Quelle grandeur se mesure en volts ?
\item Avec quel appareil ?
\end{questions}
\end{exercice}
\end{minipage}
\par\noindent\begin{minipage}{\linewidth}
\begin{exercice}[1]{Deux connexions}
\begin{questions}
\item Comment brancher un voltmètre pour mesurer la tension d’une lampe ?
\item Faut-il ouvrir sa branche pour insérer le voltmètre ?
\end{questions}
\end{exercice}
\end{minipage}
\par\noindent\begin{minipage}{\linewidth}
\begin{exercice}[1]{Un circuit ouvert}
\begin{questions}
\item Une pile non reliée à une lampe peut-elle présenter une tension ?
\item Le courant extérieur est-il alors nécessairement non nul ?
\end{questions}
\end{exercice}
\end{minipage}
\par\noindent\begin{minipage}{\linewidth}
\begin{exercice}[1]{Convertir}
\begin{questions}
\item Convertir $2{,}5\unite V$ en mV.
\item Convertir $600\unite{mV}$ en V.
\end{questions}
\end{exercice}
\end{minipage}
\par\noindent\begin{minipage}{\linewidth}
\begin{exercice}[1]{Une dérivation simple}
\pcFigureEx{figures/ex-tensions-45.pdf}{}
\begin{questions}
\item Deux lampes sont reliées aux mêmes bornes d’une source de $4{,}5\unite V$. Quelles tensions ont-elles dans le modèle ?
\item Faut-il partager 4,5 V en deux ?
\end{questions}
\end{exercice}
\end{minipage}
\clearpage\markright{Je m’entraîne}
\rubrique{Je m’entraîne}
\par\noindent\begin{minipage}{\linewidth}
\begin{exercice}[2]{Compléter une somme}
\pcFigureEx{figures/ex-tensions-serie.pdf}{}
\begin{questions}
\item En série, $U_G=9\unite V$ et $U_1=4\unite V$. Calculer $U_2$.
\item Vérifier le résultat.
\end{questions}
\end{exercice}
\end{minipage}
\par\noindent\begin{minipage}{\linewidth}
\begin{exercice}[2]{Trois récepteurs}
\begin{questions}
\item Une boucle série comporte des tensions de $1{,}2$, $2{,}0$ et $2{,}8\unite V$. Quelle tension de générateur correspond ?
\item Quelle convention est supposée ?
\end{questions}
\end{exercice}
\end{minipage}
\par\noindent\begin{minipage}{\linewidth}
\begin{exercice}[2]{Inverser les pointes}
\begin{questions}
\item Le voltmètre affiche $+3{,}0\unite V$. Que lit-on en inversant V et COM sans changer le circuit ?
\item La valeur absolue a-t-elle changé ?
\end{questions}
\end{exercice}
\end{minipage}
\par\noindent\begin{minipage}{\linewidth}
\begin{exercice}[2]{Un interrupteur idéal}
\begin{questions}
\item Quelle tension modélise-t-on aux bornes d’un interrupteur fermé idéal ?
\item Peut-il néanmoins être parcouru par un courant ?
\end{questions}
\end{exercice}
\end{minipage}
\par\noindent\begin{minipage}{\linewidth}
\begin{exercice}[2]{Même tension, autre courant}
\begin{questions}
\item Deux lampes différentes sont en dérivation. Leurs tensions sont-elles égales ?
\item Leurs intensités le sont-elles nécessairement ?
\end{questions}
\end{exercice}
\end{minipage}
\clearpage\markright{Je raisonne}
\rubrique{Je raisonne}
\par\noindent\begin{minipage}{\linewidth}
\begin{exercice}[3]{Deux lampes non identiques}
\begin{questions}
\item En série sous $6\unite V$, on mesure $1{,}8$ et $4{,}2\unite V$. Est-ce cohérent ?
\item Pourquoi une réponse de 3 V pour chacune n’était-elle pas imposée ?
\end{questions}
\end{exercice}
\end{minipage}
\par\noindent\begin{minipage}{\linewidth}
\begin{exercice}[3]{Une mesure en trop ?}
\pcFigureEx{figures/ex-voltmetre-meme-noeud.pdf}{}
\begin{questions}
\item Un élève relie les deux pointes d’un voltmètre au même nœud et lit environ $0\unite V$. Est-ce surprenant ?
\item Mesure-t-il la tension de la lampe située ailleurs ?
\end{questions}
\end{exercice}
\end{minipage}
\par\noindent\begin{minipage}{\linewidth}
\begin{exercice}[3]{Un bilan mesuré}
\begin{questions}
\item $U_G=6{,}02\unite V$, $U_1=2{,}51\unite V$, $U_2=3{,}49\unite V$. Comparer somme et source.
\item Comment conclure sans excès ?
\end{questions}
\end{exercice}
\end{minipage}
\par\noindent\begin{minipage}{\linewidth}
\begin{exercice}[3]{Une confusion dangereuse}
\pcFigureEx{figures/ex-amperemetre-danger.pdf}{Mode ampèremètre directement aux bornes de la source : montage à ne pas réaliser.}
\begin{questions}
\item Un élève veut mesurer la tension en mode ampèremètre, directement aux bornes d’une pile. Quel risque ?
\item Que doit-il vérifier avant tout branchement ?
\end{questions}
\end{exercice}
\end{minipage}
\par\noindent\begin{minipage}{\linewidth}
\begin{exercice}[3]{Retrouver le montage}
\begin{questions}
\item Deux lampes ont chacune 6 V tandis que le générateur idéal fournit 6 V. Est-ce compatible avec une série simple de récepteurs aux tensions positives ?
\item Quel montage est compatible ?
\end{questions}
\end{exercice}
\end{minipage}
\pcPage{Défi du chapitre}{Mesurer la tension électrique}
\phantomsection\label{defi-ex-15}
\pcDefiTitre{15}{Distinguer deux montages}
\begin{questions}
\item Une source idéale de $9\unite V$ alimente deux récepteurs. En montage A, $U_1=3\unite V$ et $U_2=6\unite V$. En B, les deux valent 9 V. Proposer les associations.
\item En A, peut-on déduire $I_1$ et $I_2$ séparément sans autre donnée ?
\item En B, peut-on affirmer qu’ils sont égaux ?
\end{questions}
\pcGrille
